\documentclass[11pt, a4paper]{article}

% Gemeinsame Style-Definitionen (TEXINPUTS muss templates/ enthalten — siehe scripts/build_tex.py)
% bfw-style.tex — gemeinsame Style-Definitionen für alle Handouts/Aufgabenhefte
% Voraussetzung: Kompilation mit xelatex (wegen fontspec)

\usepackage[ngerman]{babel}
\usepackage{geometry}
\geometry{a4paper, top=2.2cm, bottom=2.2cm, left=2.3cm, right=2.3cm}

\usepackage{fontspec}
% Fallback-Kette: Source Sans Pro -> Calibri -> Segoe UI -> Latin Modern Sans
% (Latin Modern ist immer mit MiKTeX vorhanden)
\IfFontExistsTF{Source Sans Pro}{%
  \setmainfont{Source Sans Pro}[Ligatures=TeX]%
}{\IfFontExistsTF{Calibri}{%
  \setmainfont{Calibri}[Ligatures=TeX]%
}{\IfFontExistsTF{Segoe UI}{%
  \setmainfont{Segoe UI}[Ligatures=TeX]%
}{%
  \setmainfont{Latin Modern Sans}[Ligatures=TeX]%
}}}
\IfFontExistsTF{Source Code Pro}{%
  \setmonofont{Source Code Pro}[Scale=0.92]%
}{\IfFontExistsTF{Consolas}{%
  \setmonofont{Consolas}[Scale=0.92]%
}{%
  \setmonofont{Latin Modern Mono}[Scale=0.92]%
}}

\usepackage{xcolor}
% Farbpalette: hell, freundlich, modern
\definecolor{bfwPrimary}{HTML}{0EA5A4}    % Petrol/Türkis - Primär
\definecolor{bfwPrimaryDark}{HTML}{0F766E} % dunkler Petrol für Headings
\definecolor{bfwAccent}{HTML}{F59E0B}     % Amber - Akzent
\definecolor{bfwSuccess}{HTML}{10B981}    % Grün - Erfolg / Lösung
\definecolor{bfwWarn}{HTML}{F97316}       % Orange - Achtung
\definecolor{bfwInk}{HTML}{0F172A}        % fast Schwarz
\definecolor{bfwInkSoft}{HTML}{475569}    % gedämpftes Grau
\definecolor{bfwBgSoft}{HTML}{F8FAFC}     % off-white
\definecolor{bfwBgTeal}{HTML}{ECFEFF}     % helles Türkis
\definecolor{bfwBgAmber}{HTML}{FEF3C7}    % helles Gelb
\definecolor{bfwBgRose}{HTML}{FFE4E6}     % helles Rosé für Eselsbrücken

\usepackage[colorlinks=true, linkcolor=bfwPrimaryDark, urlcolor=bfwPrimaryDark]{hyperref}
\usepackage{enumitem}
\setlist[itemize]{leftmargin=*, itemsep=2pt, topsep=2pt}
\setlist[enumerate]{leftmargin=*, itemsep=2pt, topsep=2pt}

\usepackage[most]{tcolorbox}
\usepackage{booktabs}
\usepackage{tikz}
\usetikzlibrary{decorations.pathmorphing, positioning, shapes.geometric, arrows.meta}

% Merkkästchen — wichtige Definition / Kernpunkt
\newtcolorbox{merksatz}[1][]{%
  enhanced, breakable,
  colback=bfwBgTeal,
  colframe=bfwPrimary,
  boxrule=0.6pt,
  arc=4pt,
  left=10pt, right=10pt, top=8pt, bottom=8pt,
  fonttitle=\bfseries\color{bfwPrimaryDark},
  title={\faLightbulb\ Merksatz},
  #1
}

% Eselsbrücke — Mnemoniks
\newtcolorbox{eselsbruecke}[1][]{%
  enhanced, breakable,
  colback=bfwBgRose,
  colframe=bfwAccent,
  boxrule=0.6pt,
  arc=4pt,
  left=10pt, right=10pt, top=8pt, bottom=8pt,
  fonttitle=\bfseries\color{bfwWarn},
  title={Eselsbrücke},
  #1
}

% Beispiel-Box
\newtcolorbox{beispiel}[1][]{%
  enhanced, breakable,
  colback=bfwBgSoft,
  colframe=bfwInkSoft,
  boxrule=0.4pt,
  arc=3pt,
  left=10pt, right=10pt, top=8pt, bottom=8pt,
  fonttitle=\bfseries\color{bfwInk},
  title={Beispiel},
  #1
}

% Lösungs-Box (für Aufgabenhefte)
\newtcolorbox{loesung}[1][]{%
  enhanced, breakable,
  colback=bfwBgAmber!40,
  colframe=bfwSuccess,
  boxrule=0.6pt,
  arc=3pt,
  left=10pt, right=10pt, top=8pt, bottom=8pt,
  fonttitle=\bfseries\color{bfwSuccess!70!black},
  title={Lösung},
  #1
}

% Glossar-Eintrag
\newcommand{\glossareintrag}[2]{%
  \par\noindent\textbf{\color{bfwPrimaryDark}#1} \enspace #2\par\smallskip
}

% Section-Stil — etwas farbiger als Standard
\usepackage{titlesec}
\titleformat{\section}
  {\Large\bfseries\color{bfwPrimaryDark}}
  {\thesection}{0.6em}{}
\titleformat{\subsection}
  {\large\bfseries\color{bfwPrimary}}
  {\thesubsection}{0.5em}{}
\titleformat{\subsubsection}
  {\normalsize\bfseries\color{bfwInk}}
  {\thesubsubsection}{0.4em}{}

% TikZ-Stil "skizzenhaft" — etwas weicher als Rough.js, aber stabil
\tikzset{
  roughbox/.style = {
    draw=bfwPrimaryDark, thick, fill=bfwBgTeal,
    rounded corners=4pt, inner sep=6pt
  },
  roughaccent/.style = {
    draw=bfwAccent, thick, fill=bfwBgAmber,
    rounded corners=4pt, inner sep=6pt
  },
  roughline/.style = {
    draw=bfwInkSoft, thick, -{Stealth[length=2.5mm]}
  }
}

% Bullet-Symbole (bewusst kein FontAwesome — vermeidet Font-Installations-Probleme)
\newcommand{\faLightbulb}{$\bullet$}
\newcommand{\faBookmark}{$\bullet$}

% Header/Footer
\usepackage{fancyhdr}
\pagestyle{fancy}
\fancyhf{}
\renewcommand{\headrulewidth}{0pt}
\fancyfoot[L]{\small\color{bfwInkSoft}\thepage}
\fancyfoot[R]{\small\color{bfwInkSoft} BFW M-064 Beschaffung im Büromanagement}


\title{\vspace*{-1.5cm}\Huge\bfseries\color{bfwPrimaryDark}Tag~8: Besprechungen nachbereiten \& Protokoll\\[0.4em]
  \large\color{bfwInkSoft}\textnormal{Vom gesprochenen Wort zur verbindlichen Niederschrift}}
\author{\textsf{Büroprozesse gestalten und Arbeitsvorgänge organisieren \textperiodcentered{} M-067}\\
\textsf{\small\copyright\ Can Siebert\,\textperiodcentered\,2026}}
\date{08.07.2026}

\begin{document}
\maketitle
\thispagestyle{empty}

\vspace{1em}
\begin{merksatz}[title={\bfseries Worum geht es heute?}]
Eine Besprechung ist erst dann erfolgreich, wenn ihre Ergebnisse \emph{festgehalten} und
\emph{umgesetzt} werden. Genau das leistet die Nachbereitung: Aufgaben werden verfolgt,
in einer Maßnahmenliste festgehalten und über einen Verteiler kommuniziert. Das wichtigste
Werkzeug dafür ist das \emph{Protokoll} --- die schriftliche Niederschrift über Inhalte,
Verlauf und Ergebnisse. In diesem Handout lernen Sie, welche Protokollarten es gibt, wie
ein Protokoll nach DIN 5008 aufgebaut ist, in welcher Sprache es geschrieben wird und
welche rechtliche Bedeutung es hat --- prüfungsrelevantes Kernwissen für die IHK-Abschlussprüfung.
\end{merksatz}

\tableofcontents
\vspace{1em}
\hrule
\vspace{1em}

\section{Besprechungen nachbereiten}

Eine durchdachte Vorbereitung ist ein wesentlicher Baustein für eine erfolgreich
verlaufende Sitzung. Doch mit dem Ende der Besprechung ist die Arbeit nicht getan ---
im Gegenteil: Erst die \emph{Nachbereitung} entscheidet darüber, ob aus Gesprächen auch
Ergebnisse werden. Wer Beschlüsse fasst, aber nicht verfolgt, verschenkt die investierte
Zeit aller Beteiligten.

\subsection{Aufgabenverfolgung und Maßnahmenliste}

Damit vereinbarte Maßnahmen nicht im Tagesgeschäft untergehen, gehört die
\emph{Aufgabenverfolgung} zur Nachbereitung. Bewährt hat sich eine \emph{To-do-} bzw.\
\emph{Maßnahmenliste}, die alle vereinbarten Aufgaben übersichtlich auflistet. Nach
DIN 5008 kann man neben dem eigentlichen Protokolltext in einer separaten Spalte auf
Termine, Aufgaben und Aktionen hinweisen; alternativ steht die Maßnahmenliste am Ende
des Inhalts.

\begin{merksatz}
Jede gute Maßnahmenliste beantwortet drei Fragen: \textbf{Was} ist zu tun (Aufgabe),
\textbf{wer} ist verantwortlich (Person) und \textbf{bis wann} (Termin)? Fehlt eine
dieser Angaben, bleibt die Aufgabe unverbindlich.
\end{merksatz}

\subsection{Verteiler}

Der \emph{Verteiler} legt fest, wer das fertige Protokoll erhält. Dazu gehören die
Teilnehmer der Besprechung, häufig aber auch entschuldigt Abwesende oder Vorgesetzte, die
über die Ergebnisse informiert sein müssen. So dient das Protokoll zugleich als
\emph{Information} für alle, die nicht (durchgängig) anwesend waren.

\section{Das Protokoll --- Begriff und Funktionen}

Das Wort \textbf{Protokoll} bedeutet \emph{Niederschrift} und meint in der kaufmännischen
Praxis eine schriftliche Aufzeichnung über Inhalte, Verlauf und Ergebnisse von
Veranstaltungen wie Sitzungen, Versammlungen, Besprechungen, Diskussionen, Konferenzen
und Verhandlungen.

Protokolle werden aus drei Gründen erstellt:

\begin{enumerate}
  \item \textbf{Information} --- z.\,B.\ für Abwesende und als Gedächtnisstütze für die Teilnehmer.
  \item \textbf{Beweis- und Dokumentationsmittel} --- Beschlüsse werden aktenkundig gemacht, Entscheidungen werden u.\,U.\ rechtswirksam.
  \item \textbf{Organisation} --- die Aufgabenverteilung wird festgehalten, Kompetenzen werden festgelegt.
\end{enumerate}

\begin{eselsbruecke}
\textbf{„I-B-O": Information --- Beweis --- Organisation.} Wer sich die drei Buchstaben
merkt, kennt die drei Funktionen jedes Protokolls.
\end{eselsbruecke}

\section{Protokollarten}

Die meisten Protokolle werden nach Stichworten geschrieben, die der Protokollführer
während der Sitzung notiert. Er hat dabei die Aufgabe, Wichtiges von Unwichtigem zu
unterscheiden --- die in der Praxis verwendeten Protokollarten sind in der Regel eine
gestraffte, strukturierte Zusammenfassung des Wesentlichen. Nach ihrer \emph{Ausführlichkeit}
(also dem Umfang) unterscheidet man:

\subsection{Wortprotokoll}

Das \textbf{Wortprotokoll} gibt alle Aussagen von Rednern oder Teilnehmern \emph{wörtlich}
wieder. Es wird z.\,B.\ bei Gerichtsverhandlungen und Parlamentsdebatten verwendet; in
der Wirtschaft findet man diese Protokollart eher selten.

\subsection{Verlaufsprotokoll}

Das \textbf{Verlaufsprotokoll} ist ausführlich und gibt den \emph{Verlauf} einer Sitzung
wieder, jedoch nicht wörtlich. Es enthält neben der Zusammenfassung des Sitzungsverlaufs
alle Vorschläge, Einwände und Argumente der Teilnehmer (bei wichtigen Beiträgen auch mit
Namenshinweis), die Aufgabenverteilung sowie die Beschlüsse im Wortlaut mit
Abstimmungsergebnissen. Es wird gewählt, wenn der Verlauf der Sitzung und das
Zustandekommen von Beschlüssen nachvollziehbar sein müssen.

\subsection{Kurzprotokoll}

Das \textbf{Kurzprotokoll} gibt die wesentlichen Inhalte und Diskussionen zusammengefasst
sowie die Beschlüsse im Wortlaut mit Abstimmungsergebnissen wieder. Es stellt
nachvollziehbar dar, wie es zu den Beschlüssen gekommen ist, und hält die Aufgabenverteilung
fest. Da Sache und Ergebnis im Vordergrund stehen, wird auf die Namensnennung der
Redebeiträge weitgehend verzichtet. Diese Protokollart eignet sich z.\,B.\ für
Abteilungsleitersitzungen, Mitarbeiterbesprechungen und Beratungsgespräche und wird in
der kaufmännischen Praxis \emph{sehr häufig} eingesetzt.

\subsection{Beschluss-/Ergebnisprotokoll}

Das \textbf{Beschluss-/Ergebnisprotokoll} enthält alle wesentlichen Entscheidungen, die
Aufgabenverteilung sowie die Beschlüsse im Wortlaut mit den Abstimmungsergebnissen.
Ergebnisprotokolle sind in der Regel sehr kurz und werden z.\,B.\ bei Projekt- und
Arbeitssitzungen eingesetzt. Für diese Protokollart eignen sich standardisierte
Protokollformulare.

\subsection{Gedächtnisprotokoll}

Das \textbf{Gedächtnisprotokoll} wird erst \emph{nach} einer Besprechung aus dem
Gedächtnis angefertigt, wenn währenddessen nicht mitgeschrieben werden konnte. Eng
verwandt ist die \emph{Akten- oder Gesprächsnotiz} --- eine deutlich kürzere
Zusammenfassung von Gesprächen, die häufig Telefongespräche dokumentiert. Beide sind
weniger formal und dienen vor allem als persönliche Gedächtnisstütze.

\subsection{Protokollarten im Vergleich}

\begin{center}
\begin{tabular}{p{3.0cm} p{4.3cm} p{4.7cm}}
\toprule
\textbf{Protokollart} & \textbf{Umfang / Inhalt} & \textbf{Typischer Einsatz}\\
\midrule
Wortprotokoll & Alle Aussagen wörtlich & Gerichtsverhandlung, Parlamentsdebatte\\
Verlaufsprotokoll & Verlauf, Argumente, Beschlüsse (nicht wörtlich) & Wenn das Zustandekommen nachvollziehbar sein muss\\
Kurzprotokoll & Wesentliches zusammengefasst, Beschlüsse im Wortlaut & Abteilungs-/Mitarbeitersitzung (sehr häufig)\\
Beschluss-/Ergebnisprotokoll & Nur Entscheidungen, Aufgaben, Beschlüsse & Projekt- und Arbeitssitzung\\
Gedächtnisprotokoll & Nachträgliche Zusammenfassung & Wenn nicht mitgeschrieben wurde\\
\bottomrule
\end{tabular}
\end{center}

\section{Protokollrahmen und -inhalt}

Nach \textbf{DIN 5008} ist ein strukturierter Protokollrahmen sinnvoll. Er gliedert sich
in einen \emph{Protokollkopf} (Protokollanfang), den \emph{Protokolltext} und einen
\emph{Protokollfuß} (Protokollschluss). Empfohlen wird, neben dem Protokolltext in einer
separaten Spalte auf Termine, Aufgaben und Aktionen hinzuweisen.

\subsection{Protokollkopf}

Im \textbf{Protokollkopf} stehen die formalen Eckdaten: Unternehmen bzw.\ Abteilung, die
Bezeichnung „Protokoll", die Besprechungsgruppe (z.\,B.\ Personalbesprechung), Datum und
Zeit, Ort, Vorsitz, Teilnehmer, Protokollführung sowie die \emph{Tagesordnung} mit den
Tagesordnungspunkten (TOP). Festzuhalten ist, welches Unternehmen (bei externen
Teilnehmern) bzw.\ welche Abteilung (bei internen Sitzungen) für die Anfertigung des
Protokolls verantwortlich ist.

\begin{merksatz}
Bei vielen Teilnehmern wird eine separate Teilnehmerliste geführt. War jemand nur
\emph{zeitweise} anwesend, ist dies zu vermerken (z.\,B.\ „Frau Dr.\ Müller bis 11:20 Uhr,
Herr Schütz ab TOP 3") --- wichtig vor allem bei Beschlüssen mit rechtlicher Relevanz.
Vermerkt wird ebenso, wenn angemeldete Teilnehmer (un)entschuldigt fernbleiben.
\end{merksatz}

\subsection{Protokolltext}

Der \textbf{Protokolltext} gibt das Wesentliche zusammengefasst und präzise wieder, ohne
zu werten. Der Protokollführer entscheidet, welche Thesen, Argumente, Fragen, Beiträge
und Stellungnahmen er als bedeutsam empfindet und berücksichtigt --- er verhält sich dabei
\emph{neutral}. Nach Erstellung wird der Inhalt vom Sitzungsvorsitzenden überprüft und
gegebenenfalls geändert. Entscheidende Kriterien sind \emph{Wichtigkeit, Vollständigkeit}
und \emph{wahrheitsgetreue Wiedergabe} der Inhalte und Ergebnisse.

\subsection{Protokollfuß}

Der \textbf{Protokollfuß} enthält das Datum der Ausfertigung, die Namen sowie ---
handschriftlich oder rechtsverbindlich elektronisch --- die Unterschriften der Leitung
(des Vorsitzenden) und des Protokollführers. Ergänzend können der \emph{Verteiler} und der
Hinweis auf \emph{Anlagen} aufgenommen werden.

\begin{beispiel}
Um das Erstellen eines Protokolls zu vereinfachen, legt man ein \emph{Protokollformular}
in Tabellenform an und nutzt für die Formularfelder die Funktionen der Textverarbeitung
(„Klicken oder tippen Sie hier, um Text einzugeben"). Werden die Tagesordnungspunkte
vorab eingegeben, lassen sich die Ergebnisse während der Sitzung direkt in das Formular
schreiben.
\end{beispiel}

\section{Sprache des Protokolls}

Der Protokolltext wird im \textbf{Präsens} (Gegenwart) geschrieben, im Ausnahmefall in
der Vergangenheit. Die Sprache ist \emph{sachlich}; der Protokollführer verzichtet auf
wertende und ausschmückende Adjektive.

\begin{center}
\begin{tabular}{p{6.0cm} p{6.0cm}}
\toprule
\textbf{Bewertende Formulierung} & \textbf{Neutrale Formulierung}\\
\midrule
Herr Bauer stellt seine überaus neue, geradezu überwältigende Produktidee vor. & Herr Bauer stellt seine Produktidee vor.\\
Das Konzept wurde leider nicht angenommen. & Das Konzept wurde nicht angenommen.\\
\bottomrule
\end{tabular}
\end{center}

\medskip
\textbf{Anträge und Beschlüsse} sind im Protokoll \emph{wörtlich} wiederzugeben --- so, wie
sie gesagt wurden. Wird ein Antrag nach der Aussprache zum Beschluss erhoben, ist der
Beschlusstext bereits vorhanden; der Antragstext fällt dann weg, es bleibt der
Beschlusstext.

\subsection{Indirekte Rede und Konjunktiv}

Wird ein Redner \emph{namentlich} erwähnt und soll sein Beitrag möglichst genau
wiedergegeben werden, so erfolgt dies in \textbf{indirekter Rede} mit dem \emph{Konjunktiv~I}
(Möglichkeitsform). Der Protokollführer wird so vom Zuhörer zum \emph{Berichterstatter}:
Er gibt die Aussage objektiv und neutral wieder, ohne die Gewähr für ihre Richtigkeit zu
übernehmen.

\begin{beispiel}
Direkte Rede: Frau Schneider: „Für die Kosten der Feier steht ein Betrag von
30.000,00 Euro zur Verfügung."\\[0.4em]
Protokoll (Konjunktiv~I): Frau Schneider teilt mit, dass für die Kosten der Feier ein
Betrag von 30.000,00 Euro zur Verfügung \emph{stehe}.
\end{beispiel}

Gibt der Protokollführer eine Aussage im \emph{Indikativ} wieder, erscheint das Gesagte
als tatsächlich und wirklich --- das widerspricht der Neutralität. Lässt sich der
Konjunktiv~I jedoch nicht vom Indikativ Präsens unterscheiden (außer beim Verb \emph{sein}
der Fall bei der 3.\ Person \emph{Plural}), wählt man als Ersatz den \textbf{Konjunktiv~II}.

\begin{center}
\begin{tabular}{l l l l l}
\toprule
\textbf{Infinitiv} & \multicolumn{2}{c}{\textbf{Präsens}} & \multicolumn{2}{c}{\textbf{Präteritum}}\\
 & Indikativ & Konjunktiv~I & Indikativ & Konjunktiv~II\\
\midrule
sein & ist & sei & war & wäre\\
werden & wird & werde & wurde & würde\\
haben & hat & habe & hatte & hätte\\
müssen & muss & müsse & musste & müsste\\
schreiben & schreibt & schreibe & schrieb & schriebe\\
gehen & geht & gehe & ging & ginge\\
\bottomrule
\end{tabular}
\end{center}

\medskip
Etwa 90~\% aller Personalformen der indirekten Rede stehen laut DUDEN in der 3.\ Person
Singular (er, sie, es) und Plural (sie). Da der Konjunktiv~II auch Zweifel am
Wahrheitsgehalt ausdrücken kann, ist er in der 3.\ Person Singular zu vermeiden und der
Konjunktiv~I einzusetzen; häufig wird der Konjunktiv~II mit \emph{würde/würden}
umschrieben.

\begin{eselsbruecke}
\textbf{„Präsens, sachlich, Konjunktiv".} Drei Wörter --- und die Sprachregeln des
Protokolls sitzen: Gegenwart als Zeitform, keine wertenden Adjektive, indirekte Rede im
Konjunktiv.
\end{eselsbruecke}

\section{Bedeutung des Protokolls}

Durch \emph{Unterschrift und Genehmigung} werden die Inhalte eines Protokolls --- Aufträge,
Bewilligungen, Vereinbarungen, Kompetenzen und Beschlüsse --- \textbf{verbindlich}. Für die
Genehmigung gibt es zwei Möglichkeiten:

\begin{itemize}
  \item Das Protokoll wird an die Teilnehmer verteilt, die eine bestimmte Frist für einen Widerspruch haben. Nach Ablauf der Frist gilt das Protokoll als genehmigt.
  \item Das Protokoll wird in der nächsten Sitzung vorgelesen, und über die Genehmigung wird abgestimmt.
\end{itemize}

\begin{merksatz}
Nur Protokolle mit der \emph{Originalunterschrift} haben eine \textbf{Beweisfunktion}. Sie
sind daher sorgfältig aufzubewahren. Zur Überprüfung der Einhaltung von Terminen,
Beschlüssen und der Erledigung von Aufgaben bieten sich Terminkalender,
Wiedervorlagemappen, Plantafeln und Computerprogramme an.
\end{merksatz}

\subsection{Moderne Praxis: digital und kollaborativ}

In der modernen Büropraxis wird das Protokoll zunehmend digital geführt. \emph{Digitale
Protokollvorlagen} (z.\,B.\ Word-Formulare mit Eingabefeldern oder Vorlagen in
Kollaborationsplattformen) sichern einen einheitlichen, DIN-gerechten Aufbau und
beschleunigen die Eingabe schon während der Sitzung. Beim \emph{kollaborativen
Protokollieren} arbeiten mehrere Personen gleichzeitig in einem gemeinsamen
Online-Dokument; Ergebnisse stehen sofort allen zur Verfügung, lassen sich versionieren
und direkt mit Aufgaben- und Terminwerkzeugen verknüpfen. Zu beachten sind Datenschutz
und Zugriffsrechte sowie --- bei rechtlicher Relevanz --- die Frage der gültigen
(elektronischen) Unterschrift.

\section{Zusammenfassung}

\begin{enumerate}
  \item \textbf{Nachbereitung} sichert die Umsetzung: Aufgabenverfolgung, To-do-/Maßnahmenliste (Was --- Wer --- Bis wann) und Verteiler.
  \item Das \textbf{Protokoll} ist eine Niederschrift mit drei Funktionen: \emph{Information, Beweis, Organisation}.
  \item Fünf \textbf{Protokollarten}: Wort-, Verlaufs-, Kurz-, Beschluss-/Ergebnis- und Gedächtnisprotokoll --- das Kurzprotokoll ist in der Praxis am häufigsten.
  \item \textbf{Protokollrahmen} nach DIN 5008: Protokollkopf (Datum/Ort/Teilnehmer/TOP), Protokolltext (sachlich, vollständig, wahrheitsgetreu) und Protokollfuß (Unterschriften, Verteiler).
  \item \textbf{Sprache}: Präsens, sachlich, indirekte Rede im Konjunktiv~I (Ersatz Konjunktiv~II); Anträge und Beschlüsse wörtlich.
  \item \textbf{Bedeutung}: Genehmigung und Unterschrift machen Inhalte verbindlich; nur Protokolle mit Originalunterschrift haben Beweisfunktion. Moderne Praxis: digitale Vorlagen und kollaboratives Protokollieren.
\end{enumerate}

\newpage

\section*{Glossar}
\addcontentsline{toc}{section}{Glossar}

\glossareintrag{Protokoll}{Schriftliche Niederschrift über Inhalte, Verlauf und Ergebnisse von Sitzungen, Besprechungen, Konferenzen und Verhandlungen.}

\glossareintrag{Aufgabenverfolgung}{Systematisches Nachhalten vereinbarter Maßnahmen nach einer Besprechung, damit Beschlüsse umgesetzt werden.}

\glossareintrag{Maßnahmenliste (To-do-Liste)}{Übersicht aller vereinbarten Aufgaben mit Angabe von Aufgabe, verantwortlicher Person und Termin.}

\glossareintrag{Verteiler}{Festlegung, wer das fertige Protokoll erhält (Teilnehmer, Abwesende, Vorgesetzte).}

\glossareintrag{Wortprotokoll}{Protokollart, die alle Aussagen wörtlich wiedergibt; typisch für Gerichtsverhandlungen und Parlamentsdebatten.}

\glossareintrag{Verlaufsprotokoll}{Ausführliche Wiedergabe des Sitzungsverlaufs mit Argumenten und Beschlüssen, jedoch nicht wörtlich.}

\glossareintrag{Kurzprotokoll}{Zusammengefasste Wiedergabe der wesentlichen Inhalte und Beschlüsse; in der kaufmännischen Praxis am häufigsten.}

\glossareintrag{Beschluss-/Ergebnisprotokoll}{Sehr kurzes Protokoll, das nur Entscheidungen, Aufgabenverteilung und Beschlüsse festhält; typisch für Projekt- und Arbeitssitzungen.}

\glossareintrag{Gedächtnisprotokoll}{Nachträglich aus dem Gedächtnis erstellte Zusammenfassung einer Besprechung.}

\glossareintrag{Protokollkopf}{Protokollanfang mit Unternehmen/Abteilung, Datum, Zeit, Ort, Vorsitz, Teilnehmern, Protokollführung und Tagesordnung.}

\glossareintrag{Protokolltext}{Sachliche, vollständige und wahrheitsgetreue Wiedergabe des Wesentlichen, geordnet nach Tagesordnungspunkten.}

\glossareintrag{Protokollfuß}{Protokollschluss mit Ausfertigungsdatum, Unterschriften von Leitung und Protokollführer, ggf.\ Verteiler und Anlagen.}

\glossareintrag{DIN 5008}{Norm für Schreib- und Gestaltungsregeln; gibt den strukturierten Protokollrahmen (Kopf, Text, Fuß) vor.}

\glossareintrag{Indirekte Rede}{Wiedergabe einer Aussage in der Möglichkeitsform; der Protokollführer wird zum Berichterstatter.}

\glossareintrag{Konjunktiv~I}{Möglichkeitsform für die indirekte Rede (z.\,B.\ „stehe", „habe"); drückt neutrale Wiedergabe ohne Gewähr aus.}

\glossareintrag{Konjunktiv~II}{Ersatzform, wenn der Konjunktiv~I nicht vom Indikativ Präsens zu unterscheiden ist (z.\,B.\ „würden", „dächten").}

\glossareintrag{Beweisfunktion}{Eigenschaft eines Protokolls mit Originalunterschrift, als Nachweis von Beschlüssen und Entscheidungen zu dienen.}

\glossareintrag{Genehmigung des Protokolls}{Verbindlichmachung durch Ablauf einer Widerspruchsfrist oder durch Abstimmung in der nächsten Sitzung.}

\glossareintrag{Kollaboratives Protokollieren}{Gemeinsames, gleichzeitiges Mitschreiben mehrerer Personen in einem digitalen Dokument.}

\end{document}
